% Source: https://artofproblemsolving.com/community/c3h1439656_efficiency_problem
% Stable identity: aops-topic-1439656-q01
Let Paul take $p$ hours, so Bob takes $p-9$ hours. In one hour Paul does $\dfrac{1}{p}$ of the job and Bob does $\dfrac{1}{p-9}$ of it, and together they do $\dfrac{1}{20}$ of it:
\[
\dfrac{1}{p}+\dfrac{1}{p-9}=\dfrac{1}{20}.
\]
Multiplying by $20p(p-9)$ gives $20(p-9)+20p=p(p-9)$, that is $40p-180=p^2-9p$, so
\[
p^2-49p+180=0,\qquad (p-45)(p-4)=0 .
\]
So the student's steps are correct. The root $p=4$ is impossible, because Bob's time $p-9$ would be negative (and also because together they could not need more time than either one alone). Hence $p=45$:
\[
\boxed{\text{Paul takes }45\text{ hours and Bob takes }36\text{ hours.}}
\]
Check: $\dfrac{1}{45}+\dfrac{1}{36}=\dfrac{4+5}{180}=\dfrac{1}{20}$.
